\section{Step 21 Propagate the subtraction borrow}\label{rsa-s21}

Subtract the modulus limb and incoming borrow from each zero-extended result limb. Keep the signed 64-bit temporary until extracting its borrow bit. Its bounded range justifies that extraction. Store the low word and move upward, then return success after either the subtracting or no-subtraction path.

Replace only the body of region \texttt{add\_borrow}, retaining its BEGIN and END comments.

\Needspace{25\baselineskip}
\MARSCodeLabel{MIPS assembly}
\begin{minted}[samepage=true]{mars-idea}
addmod_subtract:
    move $t0, $zero
    move $t1, $zero
addmod_borrow:
    daddu $t2, $a3, $t0
    lwu $t3, 0($t2)
    daddu $t4, $a2, $t0
    lwu $t4, 0($t4)
    dsubu $t3, $t3, $t4
    dsubu $t3, $t3, $t1
    dsrl $t1, $t3, 63
    sw $t3, 0($t2)
    daddiu $t0, $t0, 4
    sltiu $t2, $t0, 16
    bnez $t2, addmod_borrow
    nop
addmod_done:
    move $v0, $zero
    jr $ra
    nop
\end{minted}

Use recipe \texttt{s21-1} with the cumulative source through this edit. Stop at \texttt{function\_return} before continuing to the temporary or normal return. With the supplied caller, use \texttt{addmod128\_return} as that return breakpoint.

\subsection{Expected observation}
The carried example returns N minus 2, rather than the misleading low sum.

\Needspace{9\baselineskip}
\subsection{Check your understanding}
Why is the temporary's sign useful in this subtraction?

\begin{itemize}
  \item \textbf{A} The limb difference lies in a signed 64-bit range with no ambiguity.
  \item \textbf{B} All unsigned 64-bit differences use their sign as borrow.
  \item \textbf{C} A negative address automatically carries into the next limb.
\end{itemize}
